% !TeX program = lualatex
% !TeX encoding = utf8
% !TeX spellcheck = uk_UA
% !TeX root = ../ClassicalElectrodynamics.tex

%=========================================================
\chapter{Варіаційний принцип для рівнянь електродинаміки}\label{\currfilebase}
%=========================================================

Рівняння динаміки частинок та полів можна пов'язати з певними варіаційними принципами\index{Варіаційний принцип}, які широко використовують в теоретичній фізиці. Ці принципи пов'язані з поняттями дії та лагранжіану\index{Лагранжіан}. В цьому розділі варіаційний принцип найменшої дії\index{Принцип найменшої дії} проілюстровано на прикладі рівнянь електродинаміки. Буде подано функціонали дії, з яких можна отримати рівняння руху заряджених частинок і рівняння електромагнітного поля\index{Електромагнітне поле}. Цей розділ спирається на знання основ варіаційного числення.

%% --------------------------------------------------------
\section{Заряд у зовнішньому електромагнітному полі}
%% --------------------------------------------------------

Розглянемо траєкторію (світову лінію) точкової частинки $x^{\mu}=x^{\mu}(p)$, де $p$ --- деякий неособливий параметр. На ділянці траєкторії між точками $a=x^{\mu}(p_1)$ та $b=x^{\mu}(p_2)$ розглянемо інтеграл
\begin{equation}\label{eq:S0_free}
    S_0 = -mc\int_a^b ds
    = -mc\int_{p_1}^{p_2}
    \sqrt{\eta_{\mu\nu}\frac{dx^{\mu}}{dp}\frac{dx^{\nu}}{dp}}\,dp,
\end{equation}
що має назву \emph{дії для вільної частинки\index{Дія!вільної частинки}}; вибір коефіцієнтів забезпечує відповідність з нерелятивістським наближенням для дії. Також розглянемо інтеграл
\begin{equation}\label{eq:Sint}
    S_{\text{int}} = -\frac{q}{c}\int_a^b A_{\mu}\,dx^{\mu}
    = -\frac{q}{c}\int_{p_1}^{p_2} A_{\mu}\frac{dx^{\mu}}{dp}\,dp,
\end{equation}
що описує взаємодію частинки з електромагнітним полем; тут $A_{\mu}$ --- чотиривектор потенціалу\index{Чотиривектор!потенціал}, $q$ --- заряд частинки. Легко перевірити, що параметр $p$ можна вибрати довільно; при будь-якій заміні $p\to p'=p'(p)$ вигляд~\eqref{eq:S0_free} та~\eqref{eq:Sint} не змінюється.

Уведемо дію для частинки у зовнішньому полі:
\begin{equation*}
    S = S_0 + S_{\text{int}}.
\end{equation*}
Дія $S$ є функціоналом, що зіставляє кожній траєкторії певне число. Далі ми розглянемо першу варіацію $S$ по траєкторії частинки. Зауважимо, що повний функціонал дії є сумою дій для окремих частинок плюс дія для поля. Але при обчисленні варіацій $S$ траєкторії різних частинок вважаються незалежними, тому на цьому етапі досить розглянути дію для однієї частинки. Дія для поля не залежить від змінних частинок і не впливає на варіацію $S$ при варіаціях траєкторії частинки $x_0^{\mu}(p)\to x_0^{\mu}(p)+\delta x^{\mu}(p)$.

\emph{Варіаційний принцип найменшої дії для частинки у зовнішньому полі} можна сформулювати так: перша варіація $\delta S=\delta S_0+\delta S_{\text{int}}$ дорівнює нулю на траєкторії $x_0^{\mu}(p)$ тоді й тільки тоді, якщо ця траєкторія задовольняє рівнянням руху. Траєкторію $x_0^{\mu}(p)$ називають \emph{опорною}.

Зауважимо, що, як відомо з варіаційного числення, умова $\delta S=0$ є лише необхідною умовою мінімуму функціоналу дії\footnote{Фактично далі ми обмежуємося розглядом умови стаціонарності дії, а не її мінімуму.}. Аналіз достатніх умов мінімуму потребує додаткового розгляду.

Покажемо, що наслідком принципу найменшої дії згідно до формул~\eqref{eq:S0_free}, \eqref{eq:Sint} є релятивістські рівняння руху частинки в електромагнітному полі. Розглянемо траєкторії в околі опорної світової лінії $x_0^{\mu}(p)$. При малих змінах траєкторії $x_0^{\mu}(p)\to x_0^{\mu}(p)+\delta x^{\mu}(p)$ числове значення функціоналу $S$ змінюється. Нас цікавитимуть малі зміни дії
\begin{equation*}
    S\!\left[x_0^{\mu}+\delta x^{\mu}\right] - S\!\left[x_0^{\mu}\right]
\end{equation*}
при довільних $\delta x^{\mu}(p)$ за стандартної умови
\begin{equation}\label{eq:fixed_ends}
    \delta x^{\mu}(p_1) = \delta x^{\mu}(p_2) = 0
\end{equation}
(так звана варіація з закріпленими кінцями). Малі зміни дії розглядаємо з точністю до лінійних по $\delta x^{\mu}$ членів, тобто розгляд обмежується першою варіацією $\delta S \approx S[x_0^{\mu}+\delta x^{\mu}]-S[x_0^{\mu}]$.

Дію $S$ запишемо у вигляді
\begin{equation*}
    S = \int L\!\left(x(p),\dot{x}(p)\right)dp,
\end{equation*}
де $L(x,\dot{x}) = -mc\sqrt{\eta_{\mu\nu}\dot{x}^{\mu}\dot{x}^{\nu}} - \dfrac{q}{c}A_{\mu}(x)\dot{x}^{\mu}$; тут позначено $\dot{x}^{\alpha}\equiv dx^{\alpha}/dp$.

Покажемо, що з умови $\delta S=0$ випливають рівняння руху зарядженої частинки:
\begin{equation}\label{eq:EOM_charged}
    mc^2\frac{d^2x^{\mu}}{ds^2} = qF^{\mu}{}_{\nu}\frac{dx^{\nu}}{ds},
    \qquad\text{де}\quad
    F_{\mu\nu} = \partial_{\mu}A_{\nu} - \partial_{\nu}A_{\mu}.
\end{equation}
За стандартною процедурою варіаційного числення умова $\delta S=0$ з урахуванням~\eqref{eq:fixed_ends} приводить до \emph{рівнянь Ейлера--Лагранжа\index{Рівняння Ейлера--Лагранжа}} для функції $L(x,\dot{x})$:
\begin{equation*}
    \frac{d}{dp}\frac{\partial L}{\partial \dot{x}^{\mu}}
    = \frac{\partial L}{\partial x^{\mu}},
\end{equation*}
або, обчислюючи похідні від $L$,
\begin{equation*}
    \frac{d}{dp}\!\left(mc\,\frac{\eta_{\mu\nu}\dot{x}^{\nu}}{ds/dp}
    + \frac{q}{c}A_{\mu}\right)
    = \frac{q}{c}\frac{\partial A_{\nu}}{\partial x^{\mu}}\dot{x}^{\nu},
\end{equation*}
де $ds/dp = \sqrt{\eta_{\mu\nu}\dot{x}^{\mu}\dot{x}^{\nu}}$.

Перейдемо від параметра $p$ до $s=\int_{p_1}^{p}(ds/dp')\,dp'$:
\begin{equation*}
    \frac{d}{ds}\!\left(mc\,\eta_{\mu\nu}\frac{dx^{\nu}}{ds}
    + \frac{q}{c}A_{\mu}\right)
    = \frac{q}{c}\frac{\partial A_{\nu}}{\partial x^{\mu}}\frac{dx^{\nu}}{ds}.
\end{equation*}
Розкриваємо похідну в лівій частині:
\begin{equation*}
    mc\,\eta_{\mu\nu}\frac{d^2x^{\nu}}{ds^2}
    = \frac{q}{c}\frac{\partial A_{\nu}}{\partial x^{\mu}}\frac{dx^{\nu}}{ds}
    - \frac{q}{c}\frac{d}{ds}A_{\mu}
    = \frac{q}{c}F_{\mu\nu}\frac{dx^{\nu}}{ds}.
\end{equation*}
Результатом є рівняння, які збігаються з~\eqref{eq:EOM_charged} після підняття індексу $\mu$.

Звернемо увагу, що хоча в лагранжіані $L$ фігурує 4-вектор потенціал $A_{\mu}$, він з'являється в явному вигляді в рівняннях руху лише через тензор електромагнітного поля\index{Тензор електромагнітного поля} $F_{\alpha\beta}$, який, таким чином, і є спостережуваною величиною.

%=========================================================
\begin{problem}
Релятивістська заряджена частинка з масою $m$ та зарядом $q$ рухається у полі нерухомого заряду $Q$. Отримати рівняння для тривимірної траєкторії частинки.

\textit{Розв'язання}. Коваріантний 4-вектор-потенціал нерухомого заряду в початку координат є $A_{\mu}=(\alpha r^{-1},0,0,0)$. Якщо у виразі для дії перейти до параметру $t$, підінтегральний вираз дає тривимірну функцію Лагранжа зарядженої частинки у зовнішньому полі:
\[
    L = -mc^2\sqrt{1-\frac{v^2}{c^2}} - \frac{\alpha}{r},
    \qquad \alpha = Qq.
\]
Звідси маємо інтеграл енергії:
\begin{equation}
    \frac{mc^2}{\sqrt{1-v^2/c^2}} + \frac{\alpha}{r} = \mathcal{E}.\label{eq:energy_integral}
\end{equation}
В сферичних координатах в площині $\theta=\pi/2$ інтеграл моменту є
\begin{equation}
    \frac{mr^2\dot{\varphi}}{\sqrt{1-v^2/c^2}} = L.\label{eq:angular_integral}
\end{equation}
Комбінуючи~\eqref{eq:energy_integral} та~\eqref{eq:angular_integral}, позбавимось $\sqrt{1-v^2/c^2}$ та отримаємо:
\begin{equation*}
    r^2\dot{\varphi}\!\left(\mathcal{E} - \frac{\alpha}{r}\right) = Lc^2.
\end{equation*}
З урахуванням цього співвідношення, позначаючи $u=r^{-1}$, $u'=du/d\varphi$,
\begin{equation*}
    v^2 = \dot{r}^2 + r^2\dot{\varphi}^2
    = \frac{L^2c^4(\dot{u}'^2+u^2)}{(\mathcal{E}-\alpha u)^2}.
\end{equation*}
Звідси та з~\eqref{eq:energy_integral} маємо
\begin{equation*}
    m^2c^4 = \left(\mathcal{E}-\alpha u\right)^2
    - (u'^2+u^2)\frac{L^2c^2}{1},
\end{equation*}
або
\begin{equation*}
    u'^2 + \omega_0^2 u^2 + \frac{2\alpha\mathcal{E}}{L^2c^2}\,u
    = \frac{m^2c^4-\mathcal{E}^2}{L^2c^2},
    \qquad
    \omega_0^2 = 1 - \frac{\alpha^2}{L^2c^2}.
\end{equation*}
За $\omega_0^2>0$ рівняння має вигляд закону збереження енергії для осцилятора зі зміщеним центром; його розв'язками є тригонометричні функції. За $\omega_0^2<0$ розв'язками є гіперболічні функції; зокрема, можливе падіння на центр $r\to 0$.
\end{problem}

%% --------------------------------------------------------
\section{Дія для електромагнітного поля}
%% --------------------------------------------------------

Мета цього підрозділу --- сформулювати варіаційний принцип, з якого випливають рівняння електромагнітного поля Максвелла\index{Рівняння Максвелла}. Згідно з сучасними уявленнями, дія для системи «поле + джерела» має складатися з суми функціоналів дії для поля, для зарядів та доданка, що описує взаємодію поля із зарядами:
\begin{equation}\label{eq:S_total}
    S = S_0 + S_{\text{int}} + S_f.
\end{equation}
Оскільки поля задано в усьому просторі, функціонал дії повинен містити інтегрування по чотиривимірному об'єму, а не тільки по траєкторіях, як у випадку частинок. Тому подаватимемо $S_{\text{int}}+S_f$ у вигляді такого інтегралу:
\begin{equation}\label{eq:S_field_int}
    S_{\text{int}} + S_f
    = \int_{\Omega} d^4x\,\Lambda_{\text{tot}}\!\left(A_{\mu}(x),\,A_{\mu,\nu}(x)\right),
    \qquad \Lambda_{\text{tot}} = \Lambda_{\text{int}} + \Lambda_f.
\end{equation}
\emph{Незалежними змінними, відносно яких обчислюватиметься варіація, вважаємо компоненти 4-вектора потенціалу} $A_{\mu}$. Область інтегрування $\Omega$ може охоплювати весь простір, а може бути скінченною. Варіацію функціоналу дії розглядаємо за умови, що \emph{зміни польових функцій обертаються на нуль на межі області інтегрування} $\partial\Omega$.

\textit{Зауваження}. Якщо змінити підінтегральний вираз в~\eqref{eq:S_field_int}, додавши до нього повну 4-дивергенцію $\Lambda\to\Lambda+\partial_{\mu}f^{\mu}$, де $f$ --- будь-яка функція від польових змінних (але не від їх похідних), то при обчисленні варіації додається вираз
\begin{equation*}
    \int\partial_{\mu}\delta f^{\mu}\,d^4x = 0,
\end{equation*}
що обертається на нуль в силу граничних умов на $\partial\Omega$ після застосування формули Остроградського--Гаусса. Тому \emph{додавання повної 4-дивергенції до $\Lambda$ не впливає на рівняння, що випливають з варіаційного принципу}.

Вирази для $S_0$ та $S_{\text{int}}$ отримаємо з~\eqref{eq:S0_free}, \eqref{eq:Sint} за допомогою підсумовування дій для окремих частинок:
\begin{equation}\label{eq:S_sum}
    S_0 = -mc\int ds, \qquad S_{\text{int}} = -\frac{q}{c}\int A_{\mu}\,dx^{\mu}.
\end{equation}
Як було зазначено, ці вирази слід подати у вигляді інтегралів. Зробимо це лише для $S_{\text{int}}$, оскільки $S_0$ не впливає на виведення рівнянь поля при обчисленні варіації $\delta S$ по $\delta A_{\mu}$. Для неперервного розподілу ми зіставляємо кожному елементу заряду $dq$ його траєкторію $x_q^{\mu}(t)$; відповідно
\begin{equation*}
    dq\,\frac{dx_q^{\mu}}{dt} = \rho(\boldsymbol{x})\,d^3x\,\frac{dx_q^{\mu}}{dt}
    = j^{\mu}(\boldsymbol{x})\,d^3x,
\end{equation*}
де $\rho$ --- об'ємна густина заряду, $\boldsymbol{x}=\boldsymbol{x}_q(t)$, $j^{\mu}=\rho\,dx^{\mu}/dt$ --- 4-вектор густини струму\index{Чотиривектор!густини струму}. Тоді, переходячи в $S_{\text{int}}$ до інтегралу по області $\Omega$ чотиривимірного простору, маємо:
\begin{equation}\label{eq:Sint_4D}
    S_{\text{int}} = -\frac{1}{c}\int_{\Omega} dq\,A_{\mu}(x_e)\frac{dx_q^{\mu}}{dt}\,dt
    = -\frac{1}{c}\int_{\Omega} dt\,d^3x\,A_{\mu}(x)j^{\mu}(x)
    = -\frac{1}{c^2}\int_{\Omega} d^4x\,A_{\mu}(x)j^{\mu}(x).
\end{equation}

Виникає питання, чи не втратимо ми калібрувальну інваріантність\index{Калібрувальне перетворення} теорії завдяки тому, що в~\eqref{eq:Sint_4D} явно входить 4-потенціал? Відсутність тут якихось непорозумінь випливає в першу чергу з того, що з принципу найменшої дії буде отримано правильні рівняння поля, що є калібрувально інваріантними. Але навіть на даному етапі можна пересвідчитися, що вираз~\eqref{eq:Sint_4D} є калібрувально інваріантним, якщо врахувати закон збереження заряду\index{Закон збереження заряду}. Покажемо, що~\eqref{eq:Sint_4D} не змінюється при калібрувальних перетвореннях
\begin{equation*}
    A_{\mu} \to A_{\mu} + \partial_{\mu}\chi,
\end{equation*}
де $\chi|_{\partial\Omega}=0$ на межі області. При такому перетворенні інтеграл~\eqref{eq:Sint_4D} отримує доданок
\begin{equation*}
    -\frac{1}{c^2}\int d^4x\,\frac{\partial\chi}{\partial x^{\mu}}j^{\mu}
    = -\frac{1}{c^2}\int d^4x\,\frac{\partial}{\partial x^{\mu}}(\chi j^{\mu})
    + \frac{1}{c^2}\int d^4x\,\chi\,j^{\mu}{}_{,\mu}.
\end{equation*}
Інтеграл від дивергенції зникає на межі $\partial\Omega$, а другий доданок дорівнює нулю завдяки закону збереження заряду $j^{\mu}{}_{,\mu}=0$. Це ілюструє зв'язок калібрувальної інваріантності із збереженням заряду.

Підберемо вираз для $S_f$, що приводить до рівнянь класичної електродинаміки. Поле для пошуків $\Lambda_f$ значно звужується, якщо врахувати такі обставини. По-перше, рівняння електромагнітного поля мають однаковий вигляд в усіх інерціальних системах відліку. Тому доцільно конструювати $\Lambda_f$ як скаляр за допомогою згорток вектора $A_{\mu}$ та його похідних $A_{\mu,\alpha}$. По-друге, в рівняннях руху фігурує не безпосередньо $A_{\mu}$, а величини $F_{\mu\nu}$. Саме $F_{\mu\nu}$ мають фізичний зміст як величини, які можливо вимірювати, причому вони інваріантні відносно калібрувальних перетворень. Щоб зберегти калібрувальну інваріантність в рівняннях поля, будемо конструювати $\Lambda_f$ саме з величин $F_{\mu\nu}=\partial_{\mu}A_{\nu}-\partial_{\nu}A_{\mu}$. Ці величини задовольняють рівнянням, які запишемо у вигляді:
\begin{equation}\label{eq:Bianchi}
    \varepsilon^{\alpha\beta\gamma\delta}\cdot F_{\beta\gamma,\delta} = 0.
\end{equation}

Щоб отримати лінійні рівняння поля, природно обмежитися лише квадратичними комбінаціями $F_{\mu\nu}$, які є скалярами відносно групи Лоренца\index{Група Лоренца}. Таких комбінацій, з точністю до сталого множника, лише дві:
\begin{equation*}
    I_1 = F_{\mu\nu}F^{\mu\nu} \qquad\text{та}\qquad
    I_2 = \varepsilon^{\alpha\beta\gamma\delta}F_{\alpha\beta}F_{\gamma\delta}.
\end{equation*}
Точніше, $I_2$ --- псевдоскаляр; він змінює знак при інверсіях координат; він зводиться до повної дивергенції:
\begin{equation*}
    \varepsilon^{\alpha\beta\gamma\delta}F_{\alpha\beta}F_{\gamma\delta}
    = \varepsilon^{\alpha\beta\gamma\delta}F_{\alpha\beta}(\partial_{\gamma}A_{\delta}
    - \partial_{\delta}A_{\gamma})
    = 2\varepsilon^{\alpha\beta\gamma\delta}F_{\alpha\beta}\partial_{\gamma}A_{\delta}
    = 2\varepsilon^{\alpha\beta\gamma\delta}\partial_{\gamma}(F_{\alpha\beta}A_{\delta}),
\end{equation*}
де було враховано антисиметрію $\varepsilon^{\alpha\beta\gamma\delta}=-\varepsilon^{\alpha\beta\delta\gamma}$ та рівняння~\eqref{eq:Bianchi}, що мають місце завдяки вигляду $F_{\mu\nu}$ за побудовою. Інтеграл від $I_2$ по області $\Omega$ зводиться до інтегралу по межі $\partial\Omega$ і не дає внеску в рівняння поля всередині $\Omega$. Таким чином, в нашому розпорядженні залишається лише інваріант $I_1$, тому можна вибрати
\begin{equation}\label{eq:Sf}
    S_f = -\frac{1}{16\pi c}\int F_{\mu\nu}F^{\mu\nu}\,d^4x.
\end{equation}
Вибір сталого множника зроблено заздалегідь так, щоб забезпечити відповідність рівнянням Максвелла в гаусовій системі одиниць\index{Гаусова система одиниць}. Можна показати, що вибір іншого множника вплине лише на визначення одиниці заряду. Враховуючи внесок взаємодії поля і зарядів, покладемо
\begin{equation}\label{eq:Sf_Sint_total}
    S_f + S_{\text{int}} = \int d^4x\,\Lambda,
\end{equation}
де
\begin{equation*}
    \Lambda = -\frac{1}{16\pi c}F_{\mu\nu}F^{\mu\nu}
    - \frac{1}{c^2}A_{\mu}j^{\mu}.
\end{equation*}

Нагадаємо, що незалежними польовими змінними в~\eqref{eq:Sf_Sint_total} вважаємо $A_{\mu}$, $\mu=0,1,2,3$. Обчислимо $\delta(S_f+S_{\text{int}})=\int_{\Omega}d^4x\,\delta\Lambda$ за умови, що варіації польових змінних на межі області інтегрування зникають:
\begin{equation*}
    \delta A_{\mu}(x)\big|_{\partial\Omega} = 0.
\end{equation*}
Маємо
\begin{equation*}
    \delta S_{\text{int}} = -\frac{1}{c^2}\int_{\Omega}d^4x\,j^{\nu}(x)\,\delta A_{\nu}(x);
\end{equation*}
\begin{align*}
    \delta\int_{\Omega}d^4x\,I_1
    &= 2\int_{\Omega}d^4x\,F^{\mu\nu}\delta F_{\mu\nu}
    = 2\int_{\Omega}d^4x\,F^{\mu\nu}(\partial_{\mu}\delta A_{\nu}
    - \partial_{\nu}\delta A_{\mu}) \\
    &= 4\int_{\Omega}d^4x\,F^{\mu\nu}\partial_{\mu}\delta A_{\nu}
    = 4\int_{\Omega}d^4x\,\Big[\partial_{\mu}(F^{\mu\nu}\delta A_{\nu})
    - \partial_{\mu}F^{\mu\nu}\,\delta A_{\nu}\Big] \\
    &= -4\int_{\Omega}d^4x\,\partial_{\mu}F^{\mu\nu}\,\delta A_{\nu},
\end{align*}
де застосовано формулу Остроградського--Гаусса і враховано умову на $\partial\Omega$. Звідси
\begin{equation*}
    \delta(S_f + S_{\text{int}})
    = \int_{\Omega}d^4x\,\left(\frac{1}{4\pi c}\partial_{\mu}F^{\mu\nu}
    - \frac{1}{c^2}j^{\nu}\right)\delta A_{\nu}(x).
\end{equation*}
Згідно з принципом стаціонарної дії\index{Принцип найменшої дії}, цей вираз має дорівнювати нулю для довільних варіацій $\delta A_{\nu}(x)$ всередині $\Omega$. Це можливо лише коли вираз в дужках дорівнює нулю, або остаточно:
\begin{equation}\label{eq:Maxwell_variational}
    \partial_{\mu}F^{\mu\nu} = \frac{4\pi}{c}j^{\nu}.
\end{equation}
Разом із~\eqref{eq:Bianchi} це рівняння утворює повну систему рівнянь Максвелла для тензора електромагнітного поля $F_{\mu\nu}$.